% Experimental setup and design section for the FeRRy paper (draft, 8 Oct 2026).
% Assembled into the full paper by DeveloperDocs/paper/assemble.py, between
% methods.tex and evaluation.tex (the results). The %%TABLE:<label>%% line is
% replaced by the generated float of that label (results/exp5/paper/exp5_tables.tex).
% The study cells, arms and trial counts come from scripts/exp5/params.toml and
% the scored runs (results/exp5/scores/{b1,b2,b3}/index.md and *_arms.csv).
%%%%% BEGIN SECTION %%%%%
\section{Experimental Setup and Design}
\label{sec:setup}

\subsection{Testbed}
\label{sec:setup:testbed}

Every trial runs the full system described above: the server, each
mule, and each device are separate processes that exchange messages over TCP,
and the devices train the real intrusion-detection model on CICIoT2023.
Mission time, the radio channel, and energy come from the simulated clock and
channel described above, in the interference-limited (jittery)
regime. A trial flies four missions with a 1~MB model payload and records every
event, from which a separate scorer computes the metrics. The learned per-stop
score and the learned baseline E3 are trained in FerrySim, which also runs the
largest problems ($N=48$ and $96$).

\subsection{Cells, Budgets, and the Accuracy Target}
\label{sec:setup:cells}

A cell fixes the number of devices $N$, the number of mules $K$, and the mission
budget. Devices are static and placed uniformly at random on a $200\times200$~m
square centred on the dock, a new layout per trial; the square is the same at
every $N$, so the field grows denser as $N$ grows (FerrySim's scale cells grow
it instead). We evaluate $N\in\{6,12,24\}$ devices with one mule, and up to three
mules at $N=12$ and $18$. Each $N$ has two mission budgets, set by pilots before
the studies ran. The \emph{knee} is the smallest budget whose mean update yield
reaches 95\% of the best on a budget grid (150, 180 and 262\,s at $N=6$, 12 and
24), and the \emph{stress} budget is half of it. The accuracy target
$\tau=0.71$ is the largest value, on a 0.01 grid, that at least 80\% of the
baseline H1's trials reach at every $N$'s knee. It lies 4.3 points above the
0.667 that labelling every test flow benign would score, so time to $\tau$ measures how soon a model
becomes useful, not how good it ends; Table~\ref{tab:exp5_headline} also reports
F's and H0's final accuracy. The age cap is $S=2$.

\subsection{Arms}
\label{sec:setup:arms}

Table~\ref{tab:exp5_arms} lists the arms. \sysname{} has three, which share the
dock-time plan and differ only in how they fill the flight slot: F leaves it empty and flies the plan as
committed, FX fills it with the fixed cross-layer rule, and FQ with the learned
pair score. Every baseline is a whole scheduler,
flown in the same simulator, at the same budgets, on the same seeds, and on the
wide contact band. Each baseline's departure from its paper is stated in the
table: for instance, D3 flies our expected-connectivity variant of Cui et al.'s
index, D4 runs FedEx's tour with one transporter, and D5 is FedCS without its
resource-request phase.

%%TABLE:tab:exp5_arms%%

\paragraph{How the baselines fly}
The baselines share \sysname{}'s first stages: they cluster the devices into
stops at the wide band's range, price a stop with the same feasibility test,
and deliver in an unbudgeted, nearest-first Pass~2. Unlike \sysname{}, they
serve whole stops only, have no age cap, and use an additive deadline law: a
device's window shrinks by $5\sigma$~s after a clean contact and widens by
$10\sigma$~s after a miss, with a floor of $5\sigma$~s and no ceiling. H1 admits
stops in deadline order under the deadline and budget clauses, flies them by
priority bucket (new devices first) and then by distance from the dock, and
trims at each departure as \sysname{} does. D1--D3 rank the stops by their own
key (age of information, statistical utility, the Whittle index) and walk the
ranking under the budget clause alone, skipping what fails; D5 adds stops in
order of least marginal time while they fit one round deadline. These four fly
their admitted order, re-check only the budget at each departure, and re-admit
by their own rule. D4 flies FedEx's visit-all tour with no budget check, so it
can overrun the budget, and with three mules it splits the devices by FedEx's
own assignment rather than by sectors. E3 admits every stop and, at takeoff and
at each departure, picks the next stop by its learned Q-values among the stops
that still fit the budget, with no re-plan. It is trained in FerrySim with one
mule (learning rate $5\times10^{-4}$, $\gamma=0.99$) and flown per mule with
three.

\subsection{Metrics and Statistics}
\label{sec:setup:metrics}

The primary metric is the simulated time to $\tau$: from the first takeoff to
the first evaluation of the global model that reaches $\tau$. We also report the
network age of updates (the mean number of missions since each device's last
merged update), the round-close rate, the deadline-miss rate, the update yield
(updates merged per round), and the mission energy. Trials are paired by seed,
so every arm in a cell sees the same layout, channel and data. Each comparison
is the paired mean difference with a 95\% percentile-bootstrap CI (2000
resamples) and a Wilcoxon signed-rank $p$, Holm-corrected across each study's
comparisons. A comparison is a \emph{claim} when the CI excludes 0 and the Holm
$p<0.05$. Each study's cells, arms, primary metric and reference arm, the test
families, and the pilot rules were fixed in an analysis plan before the studies
ran.

\subsection{Experiment Design}
\label{sec:setup:design}

The evaluation answers six questions. Each is answered by one or more studies,
numbered as in our analysis plan, and each study has its own pre-registered
primary metric. Unless stated, a study flies $N=6$ devices and one mule, with
20 paired trials per cell.

\paragraph{Q1: Does \sysname{} beat the state of the art?}
Study~5.3 flies every arm of Table~\ref{tab:exp5_arms} but FQ with one and with
three mules, at the knee and stress budgets, and compares each baseline with F on
time to $\tau$. H0, which has no mule, flies one cell over the live link and is
described rather than tested. FQ joined the one-mule cells after Study~5.5's
verdict (Q5), on the same seeds; its comparisons form their own exploratory
family.

\paragraph{Q2: Where does the advantage come from?}
From 5.3's traces we split each arm's mission time into transit between stops,
dwell, and the return, upload and dock turnaround. Study~5.4 pins F to each band class at the knee budget, with
the 1~MB payload and with the measured 18.8~KB one; its primary metric is the
network age of updates. In FerrySim, an exhaustive oracle over band class,
clustering, route and per-stop band measures how far F's plan is from the best
plan.

\paragraph{Q3: How does it scale, and what does deciding cost?}
Study~5.9 grows the field to $N=12$ and $24$ with one mule and flies F, FX, H1,
D3 and D4 at each $N$'s knee and stress budgets; FQ flies the same $N=12$ cells,
on the same seeds, in Study~5.5's stack check. Study~5.11 measures (a) the
planner's wall time per plan from $N=6$ to $96$ in FerrySim, run alone on the
host; (b) one to three mules at $N=12$, and six devices per mule up to $N=18$
with three mules; and (c) FerrySim budget sweeps at $N=24$, 48 and 96, whose
knees set the budgets at which we report the share of devices served.

\paragraph{Q4: Does each part of the design matter?}
We test five claims about the design, each against an alternative that removes
or replaces one part of it, on the claim's own metric. C1, reach is a decision: F pinned to one band
class (Study~5.4). C2, one derived objective: FX without its coverage term or
its dwell term, at $N=12$, on the round-close rate (Study~5.7). C3, one deadline
in three roles: plain averaging, FedBuff~\cite{nguyen2022federated},
Async-HFL~\cite{yu2023async}, and \sysname{}'s merge with FedProx's proximal
term~\cite{li2020federated} in place of \sysname{}'s merge on the same route, on
time to $\tau$ (Study~5.1); and the round deadline of FedCS and the preferred
duration of Oort in place of the per-device deadline, on the deadline-miss rate
(Study~5.2). C4, two clocks re-decided per stop: FQ against FX in the full system
at $N=12$ (Study~5.5's stack check). C5, fairness under physical
cost: MAX-AoI, the Whittle index, and F without its coverage term or age cap,
over four missions, on the age of updates (Study~5.8). Study~5.14 also flips one
component of F at a time, with 40 trials per cell: whole stops only, the local
search only, a weighted coverage rank, the age cap raised or removed, no hover
stops, and the adaptive backhaul (F+L1) against the fixed carrier, both on the
time-varying backhaul.

\paragraph{Q5: Does learning help inside the plan?}
Study~5.5 trains FQ's learned pair score in
FerrySim at six discount factors ($\gamma\in\{0, 0.25, 0.5, 0.75, 0.9, 0.99\}$),
ten seeds each, and evaluates the median-validation seed of each on 1000
held-out episodes per cell, against FX and a one-step greedy rule. The learned
score is kept only if some $\gamma$ beats $\gamma=0$ by $\epsilon=0.01$ in
held-out return. Its stack check then flies FQ with the chosen $\gamma$ (and
with $\gamma=0$) against FX and F in the full system at $N=12$. After that
verdict, an exploratory comparison, not pre-registered, flies FQ in 5.3's
one-mule cells against E3, F, FX and D4, Holm-corrected across its eight
comparisons.

\paragraph{Q6: How robust is it?}
At the knee budget, three studies each change one condition. Study~5.12 gives each device a training time,
log-normal with median 68\,s and $\sigma=0.5$ (the median set by a pilot), and
in a second cell makes 20\% of the devices five times slower (arms F, FX, H1,
D4, D5). Study~5.13 splits the data across devices by a Dirichlet draw
($\alpha=0.1$) over the dataset's attack families (F, FX, H1, D3). Study~5.15
varies the channel one axis at a time from the jittery default: the
interference amplitude (1 and 17\,dB, against 5), the path-loss exponent (2.7,
against 2.2), and the shadowing (8\,dB, against 4) (F, FX, H1). Its cells, the
default included, fly the time-varying backhaul, so its default cell is not the main
configuration.

\paragraph{What was run}
To meet the submission date, we cut the later studies. Studies~5.1, 5.2, 5.5's
stack check, 5.7, and 5.3's two added baselines (D5 and E3) flew 20 of their
planned 40 trials per cell. We dropped 5.1's H1 route with an unbudgeted Pass~2,
5.3's relaxed budget, 5.4's range-ratio and far-device settings, 5.8's
eight-mission horizon, 5.9's two- and three-mule cells and its clean channel,
5.12's measured and 10~MB payloads, 5.13's Dirichlet $\alpha=1$ and
quantity-skew partitions, and 5.15's adaptive-backhaul arms. Study~5.6, on when
learning helps, was not run, and Study~5.10, on real radios, awaits testbed
access. The remaining trials extend the same seeds. FQ was planned in 5.3 and
5.9 but left them, as the plan specified, when Study~5.5 did not keep the
learned score; its later flights in 5.3's one-mule cells are the exploratory
addition above.
%%%%% END SECTION %%%%%
